\chapter{Solutions to Chapter 15: Matching in Observational Studies}
\label{sol:chapter15}

\begin{exercisesolution}{Linear expansions of the bias-corrected estimators}{matching residual 的计数恒等式}
本题证明 \cref{prop::matching-expansion-ate} 和 \cref{prop::linear-expansion-att}。核心不是概率论，而是一个简单的 bookkeeping：一个 unit 不仅贡献自己的 observed residual；如果它被其他 units 用作 matched unit，还会额外贡献 $K_i/M$ 份 residual。

先证明 ATE 的 expansion。令
\[
\hat R_i=Y_i-\hat\mu_{Z_i}(X_i).
\]
对 treated unit $i$，有
\begin{align*}
Y_i-M^{-1}\sum_{k\in J_i}Y_k
&=\hat\mu_1(X_i)-\hat\mu_0(X_i)
  +\{Y_i-\hat\mu_1(X_i)\}  \\
&\quad
  -M^{-1}\sum_{k\in J_i}\{Y_k-\hat\mu_0(X_k)\}
  +M^{-1}\sum_{k\in J_i}\{\hat\mu_0(X_i)-\hat\mu_0(X_k)\}.
\end{align*}
最后一项正是 treated unit 的 bias term $\hat B_i$。对 control unit $i$，同理
\begin{align*}
M^{-1}\sum_{k\in J_i}Y_k-Y_i
&=\hat\mu_1(X_i)-\hat\mu_0(X_i)
  -\{Y_i-\hat\mu_0(X_i)\}  \\
&\quad
  +M^{-1}\sum_{k\in J_i}\{Y_k-\hat\mu_1(X_k)\}
  -M^{-1}\sum_{k\in J_i}\{\hat\mu_1(X_i)-\hat\mu_1(X_k)\},
\end{align*}
最后一项也正是 control unit 的 bias term $\hat B_i$。把两种情形合并并对 $i$ 求和，得到
\begin{align*}
\hat\tau^{\textup{m}}
&=n^{-1}\sumn\{\hat\mu_1(X_i)-\hat\mu_0(X_i)\}
  +\hat B  \\
&\quad
  +n^{-1}\sumn (2Z_i-1)\hat R_i
  +n^{-1}\sumn (2Z_i-1){K_i\over M}\hat R_i.
\end{align*}
这里最后一行的第二项来自 re-indexing：若 unit $i$ 被别人匹配了 $K_i$ 次，则它的 residual 在 matched counterfactual averages 中出现 $K_i/M$ 次；treated residual 带正号，control residual 带负号。因此
\[
\hat\tau^{\textup{m}}-\hat B
=n^{-1}\sumn\left[
\hat\mu_1(X_i)-\hat\mu_0(X_i)
 +(2Z_i-1)\left(1+{K_i\over M}\right)
 \{Y_i-\hat\mu_{Z_i}(X_i)\}
\right].
\]
这就是 \cref{eq::linearATE} 中的
\[
\hat\tau^{\textup{mbc}}=n^{-1}\sumn\hat\psi_i.
\]

再证明 ATT 的 expansion。ATT 只为 treated units 填补 control potential outcome，所以
\begin{align*}
\hat\tau_\textsc{T}^{\textup{m}}
&=n_1^{-1}\sumn Z_i
\left\{
Y_i-M^{-1}\sum_{k\in J_i}Y_k
\right\}  \\
&=n_1^{-1}\sumn Z_i
\left[
Y_i-\hat\mu_0(X_i)
-M^{-1}\sum_{k\in J_i}\{Y_k-\hat\mu_0(X_k)\}
\right]
\hat B_\textsc{T}.
\end{align*}
减去 $\hat B_\textsc{T}$ 后，再把 matched control residuals 按被匹配次数重排，得到
\[
\hat\tau_\textsc{T}^{\textup{mbc}}
=n_1^{-1}\sumn
\left[
Z_i\{Y_i-\hat\mu_0(X_i)\}
-(1-Z_i){K_i\over M}\{Y_i-\hat\mu_0(X_i)\}
\right],
\]
这正是 \cref{eq::linearATT} 中的
\[
\hat\tau_\textsc{T}^{\textup{mbc}}
=n_1^{-1}\sumn \hat\psi_{\textsc{T},i}.
\]

\paragraph{Remark.}
这两个 expansions 解释了 matching variance estimator 的形式。matching with replacement 使同一个 unit 可能反复充当别人的 counterfactual；$K_i$ 正是把这种重复使用转化为 residual 权重的量。因此 matching 的不确定性不能只看原始 treated-control 差值，还必须看哪些 units 在匹配图中被重复使用。
\end{exercisesolution}

\begin{exercisesolution}{Doubly robust form of the bias-corrected matching estimator for $\tau$}{matching weights 作为非参数 propensity-score weights}
本题证明 \cref{prop::mbc-dr}。由上一题已经得到
\[
\hat\tau^{\textup{mbc}}
=n^{-1}\sumn\left[
\hat\mu_1(X_i)-\hat\mu_0(X_i)
 +(2Z_i-1)\left(1+{K_i\over M}\right)
 \{Y_i-\hat\mu_{Z_i}(X_i)\}
\right].
\]
另一方面，正文定义的 residual 为
\[
\hat R_i=
\begin{cases}
Y_i-\hat\mu_1(X_i), & Z_i=1,\\
Y_i-\hat\mu_0(X_i), & Z_i=0.
\end{cases}
\]
因此
\[
(2Z_i-1)\{Y_i-\hat\mu_{Z_i}(X_i)\}
=Z_i\hat R_i-(1-Z_i)\hat R_i.
\]
把这个恒等式代回 expansion，得到
\begin{align*}
\hat\tau^{\textup{mbc}}
&=n^{-1}\sumn\{\hat\mu_1(X_i)-\hat\mu_0(X_i)\} \\
&\quad
 +n^{-1}\sumn
\left[
\left(1+{K_i\over M}\right)Z_i\hat R_i
-\left(1+{K_i\over M}\right)(1-Z_i)\hat R_i
\right].
\end{align*}
第一行就是 outcome regression estimator
\[
\hat\tau^{\textup{reg}}
=n^{-1}\sumn\{\hat\mu_1(X_i)-\hat\mu_0(X_i)\}.
\]
于是
\[
\hat{\tau}^{\textup{mbc}}
=\hat{\tau}^{\textup{reg}}
+n^{-1}\sumn
\left\{
\left(1+{K_i\over M}\right)Z_i\hat R_i
-\left(1+{K_i\over M}\right)(1-Z_i)\hat R_i
\right\},
\]
即 \cref{prop::mbc-dr}。

\paragraph{Remark.}
这个公式把 matching 和 doubly robust estimator 放在同一张图里。DR estimator 用 $1/\hat e(X_i)$ 和 $1/\{1-\hat e(X_i)\}$ 加权 residual；bias-corrected matching estimator 用 $1+K_i/M$ 加权 residual。直觉上，一个 treated unit 若处在 propensity score 很小的区域，就会被许多 control units 选为近邻，从而 $K_i$ 大；matching 通过这种局部近邻计数来近似 inverse propensity weighting。
\end{exercisesolution}

\begin{exercisesolution}{Doubly robust form of the bias-corrected matching estimator for $\tau_\textsc{T}$}{ATT 中 $K_i/M$ 作为 odds weight}
本题证明 \cref{prop::mbc-dr-att}。由 \cref{prop::linear-expansion-att}，
\[
\hat\tau_\textsc{T}^{\textup{mbc}}
=n_1^{-1}\sumn
\left[
Z_i\{Y_i-\hat\mu_0(X_i)\}
-(1-Z_i){K_i\over M}\{Y_i-\hat\mu_0(X_i)\}
\right].
\]
正文中 ATT 的 outcome regression estimator 为
\[
\hat\tau_{\textsc{T}}^{\textup{reg}}
=n_1^{-1}\sumn Z_i\{Y_i-\hat\mu_0(X_i)\}.
\]
同时，对 control units 有
\[
\hat R_i=Y_i-\hat\mu_0(X_i),\qquad Z_i=0.
\]
因此上一式可直接写为
\[
\hat\tau_\textsc{T}^{\textup{mbc}}
=\hat\tau_{\textsc{T}}^{\textup{reg}}
-n_1^{-1}\sumn {K_i\over M}(1-Z_i)\hat R_i,
\]
这就是 \cref{prop::mbc-dr-att}。

\paragraph{Remark.}
ATT 的 DR form 只需要修正 control residual，因为 treated units 的 $Y_i$ 已经直接来自目标人群。对应地，matching 只需要用 $K_i/M$ 衡量一个 control unit 对 treated population 的代表性。这个量对应于 ATT weighting 中的 odds $e(X_i)/\{1-e(X_i)\}$，而不是 ATE weighting 中的 inverse propensity weights。
\end{exercisesolution}

\begin{exercisesolution}{Revisit Example \cref{eg::chan-bmi-ATE}}{\ri{nhanes_bmi.csv} 的 matching 分析框架}
本题要求用 matching estimator 重新分析 \cref{eg::chan-bmi-ATE} 中的 NHANES BMI 数据，并与先前章节的结果比较。目标 estimand 与 \cref{sec::ATE-application} 相同，都是参加 school meal program 对 BMI 的 average causal effect。项目 worktree 中没有找到 \ri{nhanes_bmi.csv}，因此这里不编造 matching 数值；下面给出可复现代码和应报告的比较结构。

先回顾已有结果。基于 propensity score stratification，\cref{sec::application-pscore-stratification} 给出的 point estimates 大致在
\[
-0.12 \text{ 到 } -0.27
\]
之间，standard errors 约为 $0.27$ 到 $0.28$。基于 IPW，\cref{sec::ATE-application} 报告未截断时 HT estimator 为 $-1.516$，Hajek estimator 为 $-0.156$；outcome regression 和 doubly robust estimators 分别约为 $-0.017$ 和 $-0.019$。这说明该数据中 HT 对极端 propensity-score weights 很敏感，而 regression、Hajek 和 DR routes 更接近零。

下面的代码用 \ri{Matching} package 做 $1$-$M$ matching，并检查匹配前后的 covariate balance。若把 \ri{nhanes_bmi.csv} 放在运行目录，它会输出不同 $M$ 和是否 bias adjustment 下的 matching estimates、Abadie--Imbens standard errors，以及每个 covariate 的匹配前后 standardized mean differences。

\begin{lstlisting}
library(Matching)
library(sandwich)

smd_one <- function(x, z) {
  x <- as.matrix(x)
  out <- apply(x, 2, function(v) {
    m1 <- mean(v[z == 1])
    m0 <- mean(v[z == 0])
    s1 <- var(v[z == 1])
    s0 <- var(v[z == 0])
    (m1 - m0) / sqrt((s1 + s0) / 2)
  })
  out
}

matched_smd <- function(match.obj, x) {
  xt <- x[match.obj$index.treated, , drop = FALSE]
  xc <- x[match.obj$index.control, , drop = FALSE]
  apply(xt - xc, 2, function(v) mean(v) / sd(v))
}

one_match <- function(y, z, x, M = 1, bias.adjust = TRUE) {
  fit <- Match(Y = y, Tr = z, X = x, M = M,
               BiasAdjust = bias.adjust, ties = TRUE)
  summ <- summary(fit)
  c(est = fit$est,
    se = fit$se.standard,
    ai.se = fit$se,
    n.matched = length(fit$index.treated))
}

dat <- read.csv("nhanes_bmi.csv")[, -1]
z <- dat$School_meal
y <- dat$BMI
x <- as.matrix(dat[, setdiff(names(dat), c("School_meal", "BMI"))])
x <- scale(x)

before <- smd_one(x, z)

Ms <- c(1, 2, 3, 5, 10)
tab.noadj <- t(sapply(Ms, function(m) one_match(y, z, x, M = m,
                                                 bias.adjust = FALSE)))
tab.adj <- t(sapply(Ms, function(m) one_match(y, z, x, M = m,
                                               bias.adjust = TRUE)))
rownames(tab.noadj) <- paste0("M", Ms)
rownames(tab.adj) <- paste0("M", Ms)
round(tab.noadj, 3)
round(tab.adj, 3)

fit1 <- Match(Y = y, Tr = z, X = x, M = 1,
              BiasAdjust = TRUE, ties = TRUE)
after <- matched_smd(fit1, x)
balance <- cbind(before = before, after = after)
round(balance, 3)
\end{lstlisting}

结果报告应回答三个问题。第一，bias adjustment 是否改变 point estimate；若改变明显，说明 raw matching pairs 中仍有 residual covariate imbalance。第二，$M$ 从 $1$ 增加到 $10$ 时 estimate 和 standard error 是否稳定；若 $M$ 过大导致 estimate 接近 regression route，则说明 matching 正在引入较远的 controls。第三，匹配后 standardized mean differences 是否整体缩小，特别是 \cref{fig::balance-check} 中曾经较突出的 \ri{Food_Stamp}。

\paragraph{Remark.}
这个例子的教学价值在于比较三条路线：regression route 接近零，Hajek 和 DR routes 也接近零，而 unnormalized HT 是明显 outlier。matching analysis 若设计良好，通常应更接近 Hajek/DR 而不是未截断 HT；如果 matching estimate 反而接近 HT，就要检查是否有 poor overlap、少数高杠杆 units 或匹配距离过大。
\end{exercisesolution}

\begin{exercisesolution}{Revisit Chapter \cref{sec::matching-application}}{LaLonde observational study 的多方法再分析}
本题要求重新分析 \cref{sec::matching-application} 的 LaLonde observational study，并把 outcome regression、propensity score stratification、HT、Hajek、DR 与 matching estimator 以及 experimental gold standard 比较。

正文已经给出几个锚点。实验数据中，未调整、Fisher regression adjustment 和 Lin regression adjustment 的 estimates 分别约为
\[
1794,\qquad 1676,\qquad 1622.
\]
把实验数据当作 matching study，bias-adjusted matching estimate 约为 $2119.7$，AI standard error 约为 $876.4$。在 observational counterpart \ri{cps1re74.csv} 中，简单 OLS 给出 $-8506.5$，additive OLS 给出 $1067.5$，Lin estimator 给出 $-4265.8$，说明 regression specification 非常敏感。相反，bias-adjusted $1$-$1$ matching estimate 约为 $1747.8$，AI standard error 约为 $916.6$，接近 experimental benchmark。

项目 worktree 中没有找到 \ri{cps1re74.csv}，所以下面的代码给出完整复现流程。它同时报告 regression、stratification、HT、Hajek、DR、matching，以及把同一批 estimators 应用于 matched data 后的结果。

\begin{lstlisting}
library(Matching)
library(sandwich)

hc2_se <- function(fit, name = "z") {
  sqrt(vcovHC(fit, type = "HC2")[name, name])
}

ps_fit <- function(z, x, quad = FALSE) {
  x <- as.data.frame(x)
  names(x) <- paste0("x", seq_len(ncol(x)))
  dat <- data.frame(z = z, x)
  if (!quad) {
    glm(z ~ ., data = dat, family = binomial())$fitted.values
  } else {
    glm(z ~ . + I(x1^2) + I(x2^2) + I(x7^2) + I(x8^2),
        data = dat, family = binomial())$fitted.values
  }
}

strat_est <- function(z, y, ps, K = 5) {
  qs <- quantile(ps, probs = (1:(K - 1)) / K)
  g <- cut(ps, breaks = c(-Inf, qs, Inf), labels = FALSE)
  lev <- sort(unique(g))
  est <- se2 <- numeric(length(lev))
  w <- numeric(length(lev))
  for (j in seq_along(lev)) {
    id <- g == lev[j]
    w[j] <- mean(id)
    y1 <- y[id & z == 1]
    y0 <- y[id & z == 0]
    est[j] <- mean(y1) - mean(y0)
    se2[j] <- var(y1) / length(y1) + var(y0) / length(y0)
  }
  c(est = sum(w * est), se = sqrt(sum(w^2 * se2)))
}

dr_est <- function(z, y, x, truncps = c(0, 1), quad = FALSE) {
  x <- as.data.frame(x)
  names(x) <- paste0("x", seq_len(ncol(x)))
  ps <- ps_fit(z, x, quad = quad)
  ps <- pmax(truncps[1], pmin(truncps[2], ps))

  d <- data.frame(y = y, z = z, x)
  m1 <- lm(y ~ ., data = d[d$z == 1, setdiff(names(d), "z")])
  m0 <- lm(y ~ ., data = d[d$z == 0, setdiff(names(d), "z")])
  mu1 <- predict(m1, newdata = d)
  mu0 <- predict(m0, newdata = d)

  reg <- mean(mu1 - mu0)
  ht <- mean(z * y / ps) - mean((1 - z) * y / (1 - ps))
  hajek <- sum(z * y / ps) / sum(z / ps) -
    sum((1 - z) * y / (1 - ps)) / sum((1 - z) / (1 - ps))
  dr <- reg + mean(z * (y - mu1) / ps) -
    mean((1 - z) * (y - mu0) / (1 - ps))
  c(reg = reg, HT = ht, Hajek = hajek, DR = dr)
}

dat <- read.table("cps1re74.csv", header = TRUE)
dat$u74 <- as.numeric(dat$re74 == 0)
dat$u75 <- as.numeric(dat$re75 == 0)

y <- dat$re78
z <- dat$treat
x <- as.matrix(dat[, c("age", "educ", "black",
                       "hispan", "married", "nodegree",
                       "re74", "re75", "u74", "u75")])
xs <- scale(x)

fit.unadj <- lm(y ~ z)
fit.add <- lm(y ~ z + xs)
fit.lin <- lm(y ~ z * xs)
regtab <- rbind(
  unadjusted = c(est = coef(fit.unadj)["z"],
                 se = hc2_se(fit.unadj, "z")),
  additive = c(est = coef(fit.add)["z"],
               se = hc2_se(fit.add, "z")),
  Lin = c(est = coef(fit.lin)["z"],
          se = hc2_se(fit.lin, "z"))
)
round(regtab, 2)

ps <- ps_fit(z, xs)
sapply(c(5, 10, 20), function(K) strat_est(z, y, ps, K = K))

rbind(
  notrunc = dr_est(z, y, xs, truncps = c(0, 1)),
  trunc01 = dr_est(z, y, xs, truncps = c(0.01, 0.99)),
  trunc05 = dr_est(z, y, xs, truncps = c(0.05, 0.95)),
  trunc10 = dr_est(z, y, xs, truncps = c(0.10, 0.90)),
  quad = dr_est(z, y, xs, truncps = c(0.01, 0.99), quad = TRUE)
)

matchest <- Match(Y = y, Tr = z, X = xs, BiasAdjust = TRUE)
summary(matchest)

id.m <- c(matchest$index.treated, matchest$index.control)
dr_est(z[id.m], y[id.m], xs[id.m, , drop = FALSE],
       truncps = c(0.01, 0.99))
\end{lstlisting}

报告时应以 experimental estimates $1600$--$1800$ 作为 benchmark。若 stratification、Hajek 或 DR 在合理 truncation 和 model specification 下也靠近这个区间，就说明它们与 matching 给出一致证据。若某些方法远离 benchmark，尤其是 HT 或高阶 regression specification，应进一步检查 propensity-score tail、covariate balance 和 model dependence。

\paragraph{Remark.}
LaLonde example 的要点不是证明 matching 永远优于 regression，而是展示 design stage 的重要性。matching 先把 treated 和 control units 拉到可比区域，再做 outcome analysis；这常常比在原始高度不平衡样本上直接写一个长 regression 更稳健。题目中建议把 estimators 应用于 matched data，正是为了分离 ``设计造成的改变'' 和 ``分析公式造成的改变''。
\end{exercisesolution}

\begin{exercisesolution}{Data re-analyses}{Ho, Imai, King, and Stuart 的两个 datasets}
本题要求重新分析 \citet{ho2007matching} 使用的两个 political science datasets。该论文的核心观点不是把 matching 当作最终 estimator，而是把 matching 当作 nonparametric preprocessing：先改善 covariate balance、减少 model dependence，再使用原本准备使用的 parametric outcome model。论文和 replication data 的公开页面说明，相关数据可从 Harvard Dataverse 获取；\citet{ho2011matchit} 则给出了 \ri{MatchIt} package 的软件实现。项目 worktree 中没有这些 Dataverse 数据文件，因此这里给出可复现分析模板。

对每个 dataset，建议按同一套报告结构完成：
\begin{enumerate}
\item 明确 treatment、outcome、pre-treatment covariates 和目标 estimand；
\item 在原始数据上拟合主 outcome model，记录 treatment coefficient 和 robust standard error；
\item 使用 \ri{MatchIt} 做 nearest-neighbor、optimal、full matching 或 propensity-score subclassification；
\item 用 standardized mean differences、variance ratios 和 empirical CDF differences 检查匹配前后 balance；
\item 在 matched data 上重新拟合同一个 outcome model，使用 matching weights；
\item 对不同 matching methods、calipers、ratio 和 model specifications 做敏感性比较。
\end{enumerate}

下面的 R 代码把这套流程封装成一个函数。只要把 Dataverse 下载的数据整理成包含 treatment、outcome 和 covariates 的 data frame，就可以直接复用。

\begin{lstlisting}
library(MatchIt)
library(cobalt)
library(sandwich)
library(lmtest)

analyze_ho_dataset <- function(dat, treat, outcome, covars,
                               estimand = "ATT",
                               methods = c("nearest", "optimal",
                                           "full", "subclass")) {
  f.ps <- as.formula(
    paste(treat, "~", paste(covars, collapse = " + "))
  )
  f.out <- as.formula(
    paste(outcome, "~", treat, "+", paste(covars, collapse = " + "))
  )

  fit.raw <- lm(f.out, data = dat)
  raw <- coeftest(fit.raw, vcov = vcovHC(fit.raw, type = "HC2"))[treat, ]

  out <- vector("list", length(methods))
  names(out) <- methods

  for (m in methods) {
    if (m == "nearest") {
      mat <- matchit(f.ps, data = dat, method = "nearest",
                     distance = "glm", estimand = estimand,
                     ratio = 1, caliper = 0.2)
    } else if (m == "optimal") {
      mat <- matchit(f.ps, data = dat, method = "optimal",
                     distance = "glm", estimand = estimand,
                     ratio = 1)
    } else if (m == "full") {
      mat <- matchit(f.ps, data = dat, method = "full",
                     distance = "glm", estimand = estimand)
    } else {
      mat <- matchit(f.ps, data = dat, method = "subclass",
                     distance = "glm", estimand = estimand,
                     subclass = 5)
    }

    md <- match.data(mat)
    fit <- lm(f.out, data = md, weights = weights)
    est <- coeftest(fit, vcov = vcovHC(fit, type = "HC2"))[treat, ]
    bal <- bal.tab(mat, un = TRUE, disp.v.ratio = TRUE)

    out[[m]] <- list(matchit = mat,
                     coef = est,
                     balance = bal,
                     n = nrow(md))
  }

  list(raw = raw, matched = out)
}

## Example usage after downloading and cleaning a Dataverse file:
## dat1 <- read.csv("dataset1.csv")
## res1 <- analyze_ho_dataset(
##   dat = dat1,
##   treat = "treat",
##   outcome = "y",
##   covars = c("x1", "x2", "x3", "x4")
## )
## res1$raw
## lapply(res1$matched, function(x) x$coef)
## lapply(res1$matched, function(x) x$balance)
\end{lstlisting}

最终文字报告不应只列 treatment coefficients。更重要的是展示 matching 是否减少了 model dependence：如果原始数据中不同 outcome models 给出差异很大的 treatment estimates，而 matched data 中这些 estimates 收敛到同一区间，那么这正支持 \citet{ho2007matching} 的主张。相反，如果匹配后 balance 仍差，或结论强烈依赖 caliper、ratio、subclassification 数量，就应把这种敏感性作为主要结论，而不是只报告一个最顺眼的 estimate。

\paragraph{Remark.}
这道题的重点是 design before analysis。Matching 的作用类似把 observational study 改造成一个更接近 randomized experiment 的比较；它不能替代 ignorability assumption，也不能自动解决 unobserved confounding。因此，高质量再分析必须同时报告 balance、discarded units、target population 的变化和 outcome-model sensitivity。
\end{exercisesolution}
